\begin{abstract}
Video copy localization with fixed frame descriptors compares sparse anchors and refines a
narrow band, which keeps cost sub-quadratic but makes two decisions safety-critical: whether a
weak track is a real copy that strict matching missed, and where an admitted copy starts and
ends. FastVCL answers both with discrete, fail-closed decisions from a pretrained decision
interface. A strict branch (stride-4 mutual top-1 anchors, cosine 0.70, RANSAC, banded local
dynamic programming) defines high-precision candidates. A shadow RANSAC pass over the same
anchors exposes weak tracks; a hard geometric gate filters them; a probabilistic question
estimates whether a track is a copy and admits it to the tube only above 0.75. Each segment,
from either branch, receives at most five boundary candidates (keep, left, right, both, shrink)
revalidated against local features; a multiple-choice question selects one at confidence 0.65
or higher, and otherwise the original boundary is kept. The refined segments are then merged by
a dual-axis deduplication, so one query interval that maps to two reference locations survives
as two copies. The interface never outputs coordinates or segment scores; every
output carries its source, and disabling the interface recovers the strict system. Historical
no-JEV ablations show the strict setting's recall bottleneck (R 14.36, F1 24.91, versus R 31.15,
F1 45.79 when relaxed), which motivates both decisions. The new decision ladder has not been
run: we report the design, its invariants and a falsifiable protocol, and claim no accuracy or
latency gain.
\end{abstract}
